\begin{proof}[Proof of \cref{obj:thm:weighted-separation-and-frontier}]
\begin{enumerate}
\item
\textbf{Approximation and constrained priors.}\quad
We first establish the packet lower bound. Throughout the approximation argument, assume
\[
K\ge2,\qquad 2\le M\le K^2,\qquad 0\le\epsilon\le\tfrac12,
\]
and write
\[
\mathcal E=\mathcal E_{K,M,\epsilon}^{w},
\qquad
\Delta=\Delta_{K,M,\epsilon},
\qquad
z_M(\omega)=\frac M2(1+\cos\omega),
\qquad
\vartheta_M=\arccos(2/M-1).
\]
Thus \(z_M(\vartheta_M)=1\). For integers \(N\ge2\), use
\(\mathcal J_N,k_N,G_N\) from
\cref{obj:lem:jackson-packet-localization}, and define the Lebesgue-normalized kernel
\[
J_N^{\mathrm J}(\omega)=\frac{\mathcal J_N(\omega)}{2\pi},
\qquad \omega\in\mathbb R.
\]
We will use the following two kernel moments:
\[
\int_{-\pi}^{\pi}\omega^2J_N^{\mathrm J}(\omega)\,d\omega
\le\frac{64}{N^2},
\qquad
\int_{-\pi}^{\pi}|\omega|J_N^{\mathrm J}(\omega)\,d\omega
\le\frac{32}{N}.
\]
Here is a direct derivation. Set, with removable values filled continuously,
\[
D_N(\omega)=\frac{\sin(N\omega/2)}{\sin(\omega/2)},
\qquad
\mathcal Z_N^{\mathrm J}
=\int_{-\pi}^{\pi}D_N(\omega)^4\,d\omega.
\]
The finite geometric-sum identity and orthogonality give
\[
D_N(\omega)^2
=N+2\sum_{\ell=1}^{N-1}(N-\ell)\cos(\ell\omega),
\]
\[
\int_{-\pi}^{\pi}D_N(\omega)^2\,d\omega=2\pi N,
\qquad
\mathcal Z_N^{\mathrm J}
=\frac{2\pi}{3}N(2N^2+1)
\ge\frac{4\pi}{3}N^3,
\qquad
J_N^{\mathrm J}=\frac{D_N^4}{\mathcal Z_N^{\mathrm J}}.
\]
For \(|\omega|\le\pi\), the inequality
\(|\sin(\omega/2)|\ge|\omega|/\pi\) implies
\[
\omega^2D_N(\omega)^4\le\pi^2D_N(\omega)^2.
\]
Consequently,
\[
\int_{-\pi}^{\pi}\omega^2J_N^{\mathrm J}(\omega)\,d\omega
\le
\frac{2\pi^3N}{(4\pi/3)N^3}
\le\frac{24}{N^2}
\le\frac{64}{N^2}.
\]
Finally, integrating
\[
|\omega|\le\frac N4\omega^2+\frac1N
\]
against \(J_N^{\mathrm J}\), whose integral is one, proves
\[
\int_{-\pi}^{\pi}|\omega|J_N^{\mathrm J}(\omega)\,d\omega
\le\frac N4\frac{24}{N^2}+\frac1N
\le\frac{32}{N}.
\]
% lean: Causalean.Mathlib.Analysis.JacksonApproximation.jackson_first_moment

\item
Define the nonlinear remainder and its angular version by
\[
g_\epsilon(z)
=\epsilon(1-z)_+
 +(1-2\epsilon)\frac{(z-1)_+}{1+z},
\qquad z\ge0,
\qquad
f_{M,\epsilon}(\omega)=g_\epsilon\{z_M(\omega)\},
\qquad \omega\in\mathbb R.
\]
Splitting at \(z=1\) gives
\[
\phi_\epsilon(z)=\epsilon(z-1)+g_\epsilon(z).
\]
For later integration, define the smooth curvature on the entire circle, including the kink locations, by
\[
a_{M,\epsilon}(\omega)=
\begin{cases}
\displaystyle \epsilon\frac M2\cos\omega,
&z_M(\omega)<1,\\[4pt]
\displaystyle
-\frac{4(1-2\epsilon)\{z_M'(\omega)\}^2}
       {\{1+z_M(\omega)\}^3}
+\frac{2(1-2\epsilon)z_M''(\omega)}
       {\{1+z_M(\omega)\}^2},
&z_M(\omega)>1,\\[6pt]
0,&z_M(\omega)=1.
\end{cases}
\]
On each open smooth arc this equals \(f_{M,\epsilon}''\). Since
\[
z_M'(\omega)=-\frac M2\sin\omega,
\qquad
z_M''(\omega)=-\frac M2\cos\omega,
\qquad
\{z_M'(\omega)\}^2=z_M(\omega)\{M-z_M(\omega)\},
\]
we have
\[
|a_{M,\epsilon}(\omega)|\le5M.
\]
Indeed, the lower branch is bounded by \(M/4\), while on the upper branch
\[
|a_{M,\epsilon}(\omega)|
\le
\frac{4Mz_M(\omega)}{\{1+z_M(\omega)\}^3}
+\frac{M}{\{1+z_M(\omega)\}^2}
\le5M.
\]
The scalar derivative jump is
\[
g_\epsilon'(1+)-g_\epsilon'(1-)
=\frac{1-2\epsilon}{2}+\epsilon=\frac12.
\]
Thus the angular derivative has a positive jump of size
\[
J_M^{\mathrm{jump}}=\frac{\sqrt{M-1}}2
\]
at each of \(-\vartheta_M\) and \(\vartheta_M\).

Twice integrating by parts on the three arcs
\([-\pi,-\vartheta_M]\), \([-\vartheta_M,\vartheta_M]\), and
\([\vartheta_M,\pi]\) therefore yields
\[
\begin{aligned}
\int_{-\pi}^{\pi}
f_{M,\epsilon}(\omega)k_N(\omega-\vartheta_M)\,\frac{d\omega}{2\pi}
={}&
\int_{-\pi}^{\pi}
a_{M,\epsilon}(\omega)G_N(\omega-\vartheta_M)\,
\frac{d\omega}{2\pi}\\
&+\frac{\sqrt{M-1}}{4\pi}
 \{G_N(0)+G_N(2\vartheta_M)\}.
\end{aligned}
\]
The endpoint terms at \(-\pi\) and \(\pi\) cancel by periodicity; the two remaining boundary terms are the displayed derivative jumps. The bounds in
\cref{obj:lem:jackson-packet-localization} supply universal constants
\(c_{\mathrm{loc}},C_{\mathrm{loc}}>0\) such that
\[
\frac{c_{\mathrm{loc}}}{N}\le|G_N(0)|,
\qquad
\int_{-\pi}^{\pi}|G_N(\omega)|\,\frac{d\omega}{2\pi}
\le\frac{C_{\mathrm{loc}}}{N^2}.
\]
In particular, with
\[
C_{\mathrm{sm}}=5C_{\mathrm{loc}},
\]
the smooth contribution satisfies
\[
\left|
\int_{-\pi}^{\pi}
a_{M,\epsilon}(\omega)G_N(\omega-\vartheta_M)\,
\frac{d\omega}{2\pi}
\right|
\le\frac{C_{\mathrm{sm}}M}{N^2}.
\]

To control the other kink, put
\[
\beta_M=\pi-\vartheta_M.
\]
Since \(M\ge2\), one has \(\vartheta_M\in[\pi/2,\pi)\), and hence
\[
\|2\vartheta_M\|_{\mathbb T}=2\beta_M,
\qquad
\cos\beta_M=1-\frac2M.
\]
The inequality \(1-\cos\beta_M\le\beta_M^2/2\) gives
\[
\beta_M\ge\frac2{\sqrt M},
\qquad
\|2\vartheta_M\|_{\mathbb T}\ge\frac4{\sqrt M}.
\]
Applying the localized bound for \(G_N\) gives, with
\[
C_{\mathrm{opp}}=\frac{C_{\mathrm{loc}}}{16},
\]
\[
|G_N(2\vartheta_M)|
\le
\frac{C_{\mathrm{loc}}}
     {N\{1+N\|2\vartheta_M\|_{\mathbb T}\}^2}
\le\frac{C_{\mathrm{opp}}M}{N^3}.
\]

Choose a universal integer \(A_0\) satisfying
\[
A_0>
\max\left\{
1,\frac{2C_{\mathrm{opp}}}{c_{\mathrm{loc}}},
\frac{32\pi C_{\mathrm{sm}}}{c_{\mathrm{loc}}}
\right\},
\qquad
N=A_0K.
\]
Then
\[
2C_{\mathrm{opp}}\le c_{\mathrm{loc}}A_0^2,
\qquad
32\pi C_{\mathrm{sm}}\le c_{\mathrm{loc}}A_0,
\]
so \(M\le K^2\) implies
\[
|G_N(2\vartheta_M)|\le\frac{c_{\mathrm{loc}}}{2N},
\qquad
|G_N(0)+G_N(2\vartheta_M)|
\ge\frac{c_{\mathrm{loc}}}{2N}.
\]
Also,
\[
\sqrt{M-1}\ge\frac{\sqrt M}{2},
\qquad
\frac{C_{\mathrm{sm}}M}{N^2}
\le\frac{c_{\mathrm{loc}}\sqrt M}{32\pi N}.
\]
The integration-by-parts identity consequently proves
\[
\left|
\int_{-\pi}^{\pi}
f_{M,\epsilon}(\omega)k_N(\omega-\vartheta_M)\,
\frac{d\omega}{2\pi}
\right|
\ge
\frac{c_{\mathrm{loc}}\sqrt M}{32\pi N}.
\]
% lean: CausalSmith.Stat.OptvalueVanishingoverlapRate.packetRemainder_response_lower

\item
Define the two-center packet and its normalization by
\[
H_{K,M}(\omega)
=k_N(\omega-\vartheta_M)+k_N(\omega+\vartheta_M),
\qquad
W_{K,M}
=\int_{-\pi}^{\pi}
\{1+z_M(\omega)\}|H_{K,M}(\omega)|\,\frac{d\omega}{2\pi}.
\]
The frequency exclusions in
\cref{obj:lem:jackson-packet-localization} imply
\[
\int_{-\pi}^{\pi}
P_{\mathrm{test}}\{z_M(\omega)\}H_{K,M}(\omega)\,
\frac{d\omega}{2\pi}=0,
\qquad
\deg P_{\mathrm{test}}\le K.
\]
Indeed, composing such a polynomial with \(z_M\) produces a trigonometric polynomial with frequencies of absolute value at most \(K\), whereas the translated packets have frequencies between \(2N+2\) and \(6N-2\). Translation preserves these exclusions.

Both \(f_{M,\epsilon}\) and \(k_N\) are even, so their pairings with the two translates agree. The affine term in
\(\phi_\epsilon=\epsilon(z-1)+g_\epsilon\) is annihilated. Thus
\[
\begin{aligned}
\left|
\int_{-\pi}^{\pi}
\phi_\epsilon\{z_M(\omega)\}H_{K,M}(\omega)\,
\frac{d\omega}{2\pi}
\right|
&=
2\left|
\int_{-\pi}^{\pi}
f_{M,\epsilon}(\omega)k_N(\omega-\vartheta_M)\,
\frac{d\omega}{2\pi}
\right|\\
&\ge
\frac{c_{\mathrm{loc}}}{16\pi A_0}\frac{\sqrt M}{K}.
\end{aligned}
\]

We next verify the normalization over the whole parameter range. The nonnegative Jackson kernel has a strict maximum at the periodic origin: the geometric-sum representation of \(D_N\) gives
\[
|D_N(\omega)|<N
\quad\text{when }\cos\omega<1,
\qquad
|D_N(0)|=N.
\]
Since \(\cos(2\vartheta_M)<1\),
\[
H_{K,M}(\vartheta_M)
=\mathcal J_N(0)
 +\mathcal J_N(2\vartheta_M)\cos(8N\vartheta_M)
\ge\mathcal J_N(0)-\mathcal J_N(2\vartheta_M)>0.
\]
Continuity and nonnegativity of the integrand defining \(W_{K,M}\) therefore give \(W_{K,M}>0\).

For the upper bound, use
\[
|H_{K,M}(\omega)|
\le
\mathcal J_N(\omega-\vartheta_M)
+\mathcal J_N(\omega+\vartheta_M)
\]
and the identity
\[
z_M(\omega+\vartheta_M)+z_M(\omega-\vartheta_M)
=2+(M-2)(1-\cos\omega).
\]
After translating the two periodic integrals,
\[
\begin{aligned}
W_{K,M}
&\le
\int_{-\pi}^{\pi}
\{4+(M-2)(1-\cos\omega)\}J_N^{\mathrm J}(\omega)\,d\omega\\
&\le
4+\frac{M-2}{2}
 \int_{-\pi}^{\pi}\omega^2J_N^{\mathrm J}(\omega)\,d\omega\\
&\le4+\frac{32(M-2)}{N^2}
\le36.
\end{aligned}
\]
Define
\[
c_{\mathrm p}=\frac{c_{\mathrm{loc}}}{576\pi A_0}>0.
\]
Dividing the packet response by \(W_{K,M}\le36\) proves
\[
\left|
\int_{-\pi}^{\pi}
\phi_\epsilon\{z_M(\omega)\}
\frac{H_{K,M}(\omega)}{W_{K,M}}\,
\frac{d\omega}{2\pi}
\right|
\ge c_{\mathrm p}\frac{\sqrt M}{K}.
\]
% lean: CausalSmith.Stat.OptvalueVanishingoverlapRate.packetNormalized_response_lower

\item
The packet parts in the statement are explicitly
\[
\sigma_+
=(z_M)_\#
\left[
\left(\frac{H_{K,M}}{W_{K,M}}\right)_+
\frac{d\omega}{2\pi}
\right],
\qquad
\sigma_-
=(z_M)_\#
\left[
\left(-\frac{H_{K,M}}{W_{K,M}}\right)_+
\frac{d\omega}{2\pi}
\right],
\]
where both measures inside the brackets are restricted to
\([-\pi,\pi]\). They are finite and supported on \([0,M]\), and therefore have integrable first moments. The packet is even, and the fibers of \(z_M\) consist of points with equal packet sign; these are the positive and negative Jordan parts of the pushed signed packet.

The annihilation and normalization established above give
\[
\int z^\ell\,d\sigma_+
=\int z^\ell\,d\sigma_-,
\qquad \ell=0,\ldots,K,
\]
\[
\int(1+z)\,d(\sigma_++\sigma_-)=1.
\]
Writing, as in \cref{obj:def:weighted-handle},
\[
t=\sigma_+(\mathbb R),
\qquad
u=\int z\,d\sigma_+,
\qquad
z_0=\frac{1-u}{1-t},
\]
the degree-zero and degree-one identities imply
\[
\sigma_-(\mathbb R)=t,
\qquad
\int z\,d\sigma_-=u,
\qquad
t+u=\frac12.
\]
Since \(t,u\ge0\),
\[
0\le t\le\frac12,
\qquad
z_0=\frac{1/2+t}{1-t},
\qquad
\frac12\le z_0\le2\le M.
\]
In particular \(t<1\), so the completion map is defined, and its positive-part coefficient equals \(1-t\). Set
\[
\nu_0=\sigma_-+(1-t)\delta_{z_0},
\qquad
\nu_1=\sigma_++(1-t)\delta_{z_0}.
\]
Their masses and first moments satisfy
\[
\nu_0(\mathbb R)=\nu_1(\mathbb R)=t+(1-t)=1,
\qquad
\int z\,d\nu_0=\int z\,d\nu_1=u+(1-t)z_0=1.
\]
The common atom cancels in every moment difference, so
\[
(\nu_0,\nu_1)\in\mathfrak N_{K,M},
\qquad
\mathscr H^w(\sigma_+,\sigma_-)=(\nu_0,\nu_1).
\]
It also cancels in the target difference:
\[
\begin{aligned}
\int\phi_\epsilon\,d\nu_1-\int\phi_\epsilon\,d\nu_0
&=
\int_{-\pi}^{\pi}
\phi_\epsilon\{z_M(\omega)\}
\frac{H_{K,M}(\omega)}{W_{K,M}}\,
\frac{d\omega}{2\pi}.
\end{aligned}
\]
Hence the explicit completion has gap at least
\(c_{\mathrm p}\sqrt M/K\).
% lean: CausalSmith.Stat.OptvalueVanishingoverlapRate.packetCompletion_separation_lower

\item
For any real polynomial \(P_{\mathrm{test}}\) of degree at most \(K\), define
\[
e(P_{\mathrm{test}})
=\sup_{z\in[0,M]}
\frac{|\phi_\epsilon(z)-P_{\mathrm{test}}(z)|}{1+z}.
\]
The supremum is finite by continuity. For any feasible prior pair, polynomial moment matching gives
\[
\int P_{\mathrm{test}}\,d\nu_1
=\int P_{\mathrm{test}}\,d\nu_0,
\]
and probability normalization with mean one gives
\[
\int(1+z)\,d\nu_0=\int(1+z)\,d\nu_1=2.
\]
Therefore
\[
\begin{aligned}
\left|\int\phi_\epsilon\,d\nu_1-\int\phi_\epsilon\,d\nu_0\right|
&=
\left|
\int(\phi_\epsilon-P_{\mathrm{test}})\,d\nu_1
-\int(\phi_\epsilon-P_{\mathrm{test}})\,d\nu_0
\right|\\
&\le
e(P_{\mathrm{test}})
\sum_{\iota=0}^1\int(1+z)\,d\nu_\iota
=4e(P_{\mathrm{test}}).
\end{aligned}
\]
Applying this bound to the packet completion and then taking the polynomial infimum yields, with
\[
c_{\mathrm E}=\frac{c_{\mathrm p}}4,
\]
\[
c_{\mathrm E}\frac{\sqrt M}{K}\le\mathcal E.
\]
In particular \(\mathcal E>0\).
% lean: CausalSmith.Stat.OptvalueVanishingoverlapRate.packetWeightedApproxError_lower

\item
We now prove the reverse comparison \(\mathcal E\le\Delta\) by a finite weighted dual construction. Define
\[
F_{\mathrm w}(z)=\frac{\phi_\epsilon(z)}{1+z},
\qquad
\mathcal V_K^{\mathrm w}
=\left\{
z\mapsto\frac{P_{\mathrm{test}}(z)}{1+z}:
\deg P_{\mathrm{test}}\le K
\right\}
\subset C([0,M],\mathbb R).
\]
This is a finite-dimensional subspace. Its closed ball of radius
\[
R_{\mathrm{best}}=2\|F_{\mathrm w}\|_\infty+1
\]
is compact, so the distance from \(F_{\mathrm w}\) attains a minimum there. The zero function gives distance at most \(\|F_{\mathrm w}\|_\infty\), while a function outside this ball has distance at least
\[
\|v_{\mathrm w}\|_\infty-\|F_{\mathrm w}\|_\infty
>\|F_{\mathrm w}\|_\infty+1.
\]
Thus a polynomial \(P_{\mathrm{best}}\) attains the global infimum:
\[
\deg P_{\mathrm{best}}\le K,
\qquad
e(P_{\mathrm{best}})=\mathcal E.
\]
Define its weighted residual and extremal set by
\[
r_{\mathrm{best}}(z)
=\frac{\phi_\epsilon(z)-P_{\mathrm{best}}(z)}{1+z},
\qquad
\mathcal A_{\mathrm{ext}}
=\{z\in[0,M]:|r_{\mathrm{best}}(z)|=\mathcal E\}.
\]

There exist ordered nodes and an orientation satisfying
\[
0\le z_0^{\mathrm{dual}}<\cdots<z_{K+1}^{\mathrm{dual}}\le M,
\qquad
o_{\mathrm{dual}}\in\{-1,1\},
\]
\[
r_{\mathrm{best}}(z_\ell^{\mathrm{dual}})
=o_{\mathrm{dual}}(-1)^\ell\mathcal E,
\qquad \ell=0,\ldots,K+1.
\]
To see this, suppose such nodes are absent. The positive and negative extremal sets are disjoint compact sets. If one is empty, the constant polynomial with the other sign agrees with the residual sign on all extrema. If both are nonempty, their positive separation permits a finite ordered cover by sufficiently small neighborhoods. Consecutive sign blocks in that cover have at most \(K\) transitions: \(K+1\) transitions would furnish the forbidden \(K+2\) alternating nodes. Placing one root in each transition gap and choosing the overall sign produces a polynomial \(Q_{\mathrm{sign}}\) such that
\[
\deg Q_{\mathrm{sign}}\le K,
\qquad
r_{\mathrm{best}}(z)Q_{\mathrm{sign}}(z)>0
\quad(z\in\mathcal A_{\mathrm{ext}}).
\]
Put
\[
v_{\mathrm{sign}}(z)=\frac{Q_{\mathrm{sign}}(z)}{1+z}.
\]
Define the set where the perturbation has the wrong sign or vanishes by
\[
\mathcal B_{\mathrm{bad}}
=\{z\in[0,M]:r_{\mathrm{best}}(z)v_{\mathrm{sign}}(z)\le0\}.
\]
This set is compact. If it is nonempty, the continuous absolute residual attains its maximum there, and
\[
\max_{z\in\mathcal B_{\mathrm{bad}}}|r_{\mathrm{best}}(z)|<\mathcal E.
\]
Indeed, equality at a maximizing point would make that point an extremum, where the residual and perturbation have strictly positive product. Define the threshold in both cases by
\[
b_{\mathrm{ext}}=
\begin{cases}
\displaystyle\frac{\max_{z\in\mathcal B_{\mathrm{bad}}}|r_{\mathrm{best}}(z)|+\mathcal E}{2},
&\mathcal B_{\mathrm{bad}}\ne\varnothing,\\[6pt]
\mathcal E/2,&\mathcal B_{\mathrm{bad}}=\varnothing.
\end{cases}
\]
Thus \(b_{\mathrm{ext}}<\mathcal E\), and in either case
\[
|r_{\mathrm{best}}(z)|\ge b_{\mathrm{ext}}
\quad\Longrightarrow\quad
r_{\mathrm{best}}(z)v_{\mathrm{sign}}(z)>0.
\]
The high-residual set
\[
\mathcal A_{\mathrm{high}}
=\{z\in[0,M]:|r_{\mathrm{best}}(z)|\ge b_{\mathrm{ext}}\}
\]
is compact and nonempty, since the residual attains its norm \(\mathcal E\). Consequently the following extrema exist:
\[
m_{\mathrm{ext}}
=\min_{z\in\mathcal A_{\mathrm{high}}}
 r_{\mathrm{best}}(z)v_{\mathrm{sign}}(z)>0,
\qquad
B_{\mathrm{sign}}
=\max_{z\in[0,M]}|v_{\mathrm{sign}}(z)|>0.
\]
The second constant is positive because the perturbation is nonzero at every residual extremum. These definitions give
\[
r_{\mathrm{best}}(z)v_{\mathrm{sign}}(z)\ge m_{\mathrm{ext}}
\quad\text{if }|r_{\mathrm{best}}(z)|\ge b_{\mathrm{ext}},
\qquad
|v_{\mathrm{sign}}(z)|\le B_{\mathrm{sign}}
\quad(z\in[0,M]).
\]
Choose
\[
\eta_{\mathrm{pert}}
=\min\left\{
\frac{m_{\mathrm{ext}}}{B_{\mathrm{sign}}^2},
\frac{\mathcal E-b_{\mathrm{ext}}}{2B_{\mathrm{sign}}}
\right\}>0.
\]
On the first region,
\[
\begin{aligned}
\{r_{\mathrm{best}}-\eta_{\mathrm{pert}}v_{\mathrm{sign}}\}^2
&\le
\mathcal E^2-2\eta_{\mathrm{pert}}m_{\mathrm{ext}}
+\eta_{\mathrm{pert}}^2B_{\mathrm{sign}}^2\\
&\le\mathcal E^2-\eta_{\mathrm{pert}}m_{\mathrm{ext}}
<\mathcal E^2.
\end{aligned}
\]
On its complement,
\[
|r_{\mathrm{best}}-\eta_{\mathrm{pert}}v_{\mathrm{sign}}|
<
b_{\mathrm{ext}}+\frac{\mathcal E-b_{\mathrm{ext}}}{2}
<\mathcal E.
\]
Continuity on the compact interval makes the supremum strictly smaller than \(\mathcal E\), contradicting optimality of
\(P_{\mathrm{best}}+\eta_{\mathrm{pert}}Q_{\mathrm{sign}}\). This proves the alternating-node assertion.
% lean: CausalSmith.Stat.OptvalueVanishingoverlapRate.weightedApproxError_alternating_witness

\item
For these nodes, define the Lagrange weights and their two normalizations by
\[
\lambda_\ell^{\mathrm{Lag}}
=\prod_{\substack{0\le r\le K+1\\r\ne\ell}}
(z_\ell^{\mathrm{dual}}-z_r^{\mathrm{dual}})^{-1},
\qquad
\lambda_\ell^{\mathrm{norm}}
=\frac{\lambda_\ell^{\mathrm{Lag}}}
       {\sum_{r=0}^{K+1}|\lambda_r^{\mathrm{Lag}}|},
\]
\[
D_{\mathrm{dual}}
=\sum_{\ell=0}^{K+1}
|\lambda_\ell^{\mathrm{norm}}|(1+z_\ell^{\mathrm{dual}}),
\qquad
\lambda_\ell^{\mathrm{dual}}
=\frac{\lambda_\ell^{\mathrm{norm}}}{D_{\mathrm{dual}}}.
\]
The node differences are nonzero, and the two normalization denominators are positive; moreover, \(D_{\mathrm{dual}}\ge1\).
Comparing the coefficient of degree \(K+1\) in the Lagrange interpolation formula gives
\[
\sum_{\ell=0}^{K+1}
\lambda_\ell^{\mathrm{norm}}P_{\mathrm{test}}(z_\ell^{\mathrm{dual}})
=0
\qquad(\deg P_{\mathrm{test}}\le K).
\]
The ordered nodes also give
\[
\lambda_\ell^{\mathrm{norm}}
=(-1)^{K+1-\ell}|\lambda_\ell^{\mathrm{norm}}|.
\]
Combining this sign with the alternating residual yields
\[
\begin{aligned}
\sum_{\ell=0}^{K+1}
\lambda_\ell^{\mathrm{norm}}\phi_\epsilon(z_\ell^{\mathrm{dual}})
&=
\sum_{\ell=0}^{K+1}
\lambda_\ell^{\mathrm{norm}}
\{\phi_\epsilon(z_\ell^{\mathrm{dual}})
-P_{\mathrm{best}}(z_\ell^{\mathrm{dual}})\}\\
&=o_{\mathrm{dual}}(-1)^{K+1}\mathcal E D_{\mathrm{dual}}.
\end{aligned}
\]
Consequently,
\[
\sum_{\ell=0}^{K+1}
\lambda_\ell^{\mathrm{dual}}(z_\ell^{\mathrm{dual}})^r=0
\quad(r=0,\ldots,K),
\]
\[
\sum_{\ell=0}^{K+1}
|\lambda_\ell^{\mathrm{dual}}|(1+z_\ell^{\mathrm{dual}})=1,
\qquad
\left|
\sum_{\ell=0}^{K+1}
\lambda_\ell^{\mathrm{dual}}\phi_\epsilon(z_\ell^{\mathrm{dual}})
\right|=\mathcal E.
\]
Split the weights into positive and negative parts:
\[
\sigma_+^{\mathrm{dual}}
=\sum_{\ell=0}^{K+1}(\lambda_\ell^{\mathrm{dual}})_+
 \delta_{z_\ell^{\mathrm{dual}}},
\qquad
\sigma_-^{\mathrm{dual}}
=\sum_{\ell=0}^{K+1}(-\lambda_\ell^{\mathrm{dual}})_+
 \delta_{z_\ell^{\mathrm{dual}}}.
\]
Define
\[
t_{\mathrm{dual}}=\sigma_+^{\mathrm{dual}}(\mathbb R),
\qquad
u_{\mathrm{dual}}=\int z\,d\sigma_+^{\mathrm{dual}},
\qquad
z_{\mathrm{dual}}^{\mathrm{comp}}
=\frac{1-u_{\mathrm{dual}}}{1-t_{\mathrm{dual}}}.
\]
The degree-zero and degree-one identities and the weighted normalization give
\[
t_{\mathrm{dual}}+u_{\mathrm{dual}}=\frac12,
\qquad
0\le t_{\mathrm{dual}}\le\frac12,
\qquad
\frac12\le z_{\mathrm{dual}}^{\mathrm{comp}}\le2.
\]
The same mass, mean, and moment calculation used for the packet completion shows that
\[
\nu_0^{\mathrm{dual}}
=\sigma_-^{\mathrm{dual}}
 +(1-t_{\mathrm{dual}})\delta_{z_{\mathrm{dual}}^{\mathrm{comp}}},
\qquad
\nu_1^{\mathrm{dual}}
=\sigma_+^{\mathrm{dual}}
 +(1-t_{\mathrm{dual}})\delta_{z_{\mathrm{dual}}^{\mathrm{comp}}}
\]
form a pair in \(\mathfrak N_{K,M}\), with target gap \(\mathcal E\). Hence
\[
\mathcal E\le\Delta.
\]
The feasible set is nonempty, since \((\delta_1,\delta_1)\) is feasible. Its gaps are bounded by \(4e(0)\), by the polynomial residual bound above. Taking the polynomial infimum in that bound, followed by the feasible-pair supremum, therefore gives
\[
\Delta\le4\mathcal E.
\]
% lean: CausalSmith.Stat.OptvalueVanishingoverlapRate.weightedApproxError_le_constrainedPriorSeparation

\item
For the approximation upper bound, the square-root coordinate gives the required uniform angular regularity. Define, for \(\xi\ge0\),
\[
b_{\mathrm{lo}}(\xi)=(1-\xi^2)_+,
\qquad
b_{\mathrm{hi}}(\xi)=\frac{(\xi^2-1)_+}{1+\xi^2}.
\]
For \(0\le\xi\le\eta\), splitting at \(1\) proves
\[
|b_{\mathrm{lo}}(\xi)-b_{\mathrm{lo}}(\eta)|
\le2(\eta-\xi),
\qquad
|b_{\mathrm{hi}}(\xi)-b_{\mathrm{hi}}(\eta)|
\le2(\eta-\xi).
\]
For the lower branch, when \(\eta\le1\), the difference equals
\((\eta-\xi)(\eta+\xi)\); when \(\xi\le1\le\eta\), it equals
\(1-\xi^2\le2(1-\xi)\); when \(1\le\xi\), it vanishes.
For the upper branch, when \(1\le\xi\),
\[
b_{\mathrm{hi}}(\eta)-b_{\mathrm{hi}}(\xi)
=\frac{2(\eta-\xi)(\eta+\xi)}
       {(1+\xi^2)(1+\eta^2)}
\le2(\eta-\xi).
\]
When \(\xi\le1\le\eta\), the bound follows from
\[
\frac{\eta^2-1}{1+\eta^2}\le2(\eta-1)\le2(\eta-\xi);
\]
when \(\eta\le1\), the difference vanishes. Interchanging \(\xi,\eta\) covers the opposite ordering. It follows that
\[
|g_\epsilon(\xi^2)-g_\epsilon(\eta^2)|
\le
2\{\epsilon+(1-2\epsilon)\}|\xi-\eta|
\le2|\xi-\eta|.
\]
Using the half-angle identity
\[
z_M(\omega)=\{\sqrt M\,|\cos(\omega/2)|\}^2
\]
and the Lipschitz bound for cosine gives
\[
|f_{M,\epsilon}(\omega)-f_{M,\epsilon}(\upsilon)|
\le\sqrt M\,|\omega-\upsilon|,
\qquad \omega,\upsilon\in\mathbb R.
\]

Set
\[
N_{\mathrm{up}}=\lfloor K/2\rfloor+1,
\qquad
T_{\mathrm{up}}(\omega)
=\int_{-\pi}^{\pi}
f_{M,\epsilon}(\omega+\upsilon)
J_{N_{\mathrm{up}}}^{\mathrm J}(\upsilon)\,d\upsilon.
\]
Convolution with the even Jackson kernel makes \(T_{\mathrm{up}}\) an even trigonometric polynomial of degree at most
\(2(N_{\mathrm{up}}-1)\le K\). Positivity, normalization, and the first kernel moment give
\[
|f_{M,\epsilon}(\omega)-T_{\mathrm{up}}(\omega)|
\le
\sqrt M\int_{-\pi}^{\pi}
|\upsilon|J_{N_{\mathrm{up}}}^{\mathrm J}(\upsilon)\,d\upsilon
\le\frac{32\sqrt M}{N_{\mathrm{up}}}.
\]
Writing its cosine modes as Chebyshev polynomials supplies a real polynomial \(Q_{\mathrm{up}}\) such that
\[
T_{\mathrm{up}}(\omega)=Q_{\mathrm{up}}(\cos\omega),
\qquad
\deg Q_{\mathrm{up}}\le2(N_{\mathrm{up}}-1).
\]
Define
\[
P_{\mathrm{up}}(z)
=\epsilon(z-1)+Q_{\mathrm{up}}(2z/M-1).
\]
Then \(\deg P_{\mathrm{up}}\le K\), and every \(z\in[0,M]\) can be written as
\(z_M(\omega)\). Since \(K\le2N_{\mathrm{up}}\),
\[
|\phi_\epsilon(z)-P_{\mathrm{up}}(z)|
\le\frac{32\sqrt M}{N_{\mathrm{up}}}
\le\frac{64\sqrt M}{K}.
\]
Division by \(1+z\ge1\), followed by the polynomial infimum, proves
\[
\mathcal E\le\frac{64\sqrt M}{K}.
\]
% lean: CausalSmith.Stat.OptvalueVanishingoverlapRate.weightedApproxError_upper

\item
Choose the universal constants
\[
c=\min\{\min(c_{\mathrm E},c_{\mathrm p}),1\},
\qquad
C=257.
\]
They satisfy \(0<c<C\), and the preceding bounds yield
\[
c\frac{\sqrt M}{K}
\le\mathcal E
\le\Delta
\le4\mathcal E
\le256\frac{\sqrt M}{K}
\le C\frac{\sqrt M}{K}.
\]
The packet completion retains its gap because \(c\le c_{\mathrm p}\).
All the packet estimates used the closed range
\(0\le\epsilon\le1/2\) and the non-strict support bound \(M\le K^2\), so they include both overlap endpoints and \(M=K^2\).
% lean: CausalSmith.Stat.OptvalueVanishingoverlapRate.weighted_approximation_and_packet

\item
\textbf{Sparse comparison.}\quad
We next establish an observed minimax lower bound for all alphabets. For this part, assume
\[
n\ge1,\qquad d\ge2,\qquad 0<\epsilon\le\tfrac12,
\]
and define
\[
L_d=\log(ed),
\qquad
r_{n,d,\epsilon}=\frac{d}{n\epsilon L_d},
\qquad
\rho_{n,d,\epsilon}=\min\{1,r_{n,d,\epsilon}\}.
\]
In particular \(L_d=1+\log d\ge1\) and \(\rho_{n,d,\epsilon}\ge0\).

Here is the likelihood calibration used in both large-alphabet constructions. For integers \(K_{\mathrm{test}}\ge1\) and real \(\tau\ge0\), define
\[
\mathcal T_{K_{\mathrm{test}}}(\tau)
=\sum_{\ell>K_{\mathrm{test}}}\frac{\tau^\ell}{\ell!},
\qquad
\beta_{\mathrm{tail}}=\frac1{16e}.
\]
When \(\tau\le\beta_{\mathrm{tail}}K_{\mathrm{test}}\), the factorial bound
\(\ell!\ge(\ell/e)^\ell\) gives, for \(\ell>K_{\mathrm{test}}\),
\[
\frac{\tau^\ell}{\ell!}
\le\left(\frac{e\tau}{\ell}\right)^\ell
\le4^{-\ell}.
\]
Summing the geometric tail therefore gives
\[
\mathcal T_{K_{\mathrm{test}}}(\tau)
\le4^{-K_{\mathrm{test}}},
\qquad
\sqrt{\mathcal T_{K_{\mathrm{test}}}(\tau)}
\le2^{-K_{\mathrm{test}}}.
\]
Define
\[
C_{\mathrm{tv}}=\frac{|\log8|+2}{\log2}.
\]
If \(K_{\mathrm{test}}\ge C_{\mathrm{tv}}L_d\), then
\[
K_{\mathrm{test}}\log2
\ge(|\log8|+2)(1+\log d)
\ge\log(8d),
\]
and consequently
\[
d\sqrt{\mathcal T_{K_{\mathrm{test}}}(\tau)}
\le d\,2^{-K_{\mathrm{test}}}\le\frac18.
\]
% lean: CausalSmith.Stat.OptvalueVanishingoverlapRate.sparse_geometric_degree_calibration

\item
Consider first the sparse experiment of
\cref{obj:lem:sparse-likelihood}, with inflation \(B>2\) and a constrained pair of degree \(K_{\mathrm{test}}\). Under prior \(\nu_\iota\), \(\iota\in\{0,1\}\), draw cell intensities independently with that prior, and observe all four independent Poisson counts in each cell.

For one cell, let \(\mathbb Q_\star\) be the reference count law at \(z_\star=M/2\), and define the mixture densities
\[
h_{\mathrm{sp},\iota}
=\int L_z\,d\nu_\iota(z),
\qquad \iota\in\{0,1\}.
\]
The likelihood ratios \(L_z\) and Gram coefficient \(\alpha\) are those in
\cref{obj:lem:sparse-likelihood}; at a zero alternative mean, the count factor uses \(0^0=1\). Its Gram identity, followed by expansion of the exponential, gives
\[
\begin{aligned}
\int(h_{\mathrm{sp},1}-h_{\mathrm{sp},0})^2\,d\mathbb Q_\star
&=
\sum_{\ell>K_{\mathrm{test}}}\frac{\alpha^\ell}{\ell!}
\left\{
\int(z-z_\star)^\ell\,d(\nu_1-\nu_0)(z)
\right\}^2\\
&\le
4\mathcal T_{K_{\mathrm{test}}}(\tau_{\mathrm{sp}}),
\qquad
\tau_{\mathrm{sp}}=\frac{\alpha M^2}{4}.
\end{aligned}
\]
The terms through degree \(K_{\mathrm{test}}\) vanish by moment matching, and the remaining moments have absolute value at most
\(2(M/2)^\ell\). Thus Cauchy--Schwarz and telescoping products imply
\[
\operatorname{TV}
\left(\mathbb Q_{\mathrm{sp},0}^{\otimes d},
      \mathbb Q_{\mathrm{sp},1}^{\otimes d}\right)
\le
d\sqrt{\mathcal T_{K_{\mathrm{test}}}(\tau_{\mathrm{sp}})},
\]
where
\[
\mathbb Q_{\mathrm{sp},\iota}
=h_{\mathrm{sp},\iota}\,\mathbb Q_\star.
\]
Set
\[
\eta_{\mathrm{sp}}=2\beta_{\mathrm{tail}}.
\]
The bound on \(\alpha\) in
\cref{obj:lem:sparse-likelihood} gives
\[
\tau_{\mathrm{sp}}\le\frac{Bn\epsilon M}{2d}.
\]
Hence
\[
M\le\frac{\eta_{\mathrm{sp}}dK_{\mathrm{test}}}{Bn\epsilon},
\qquad
K_{\mathrm{test}}\ge C_{\mathrm{tv}}L_d
\]
ensure that the full product-mixture total variation is at most \(1/8\).
% lean: CausalSmith.Stat.OptvalueVanishingoverlapRate.sparse_product_mixture_tv_le_one_eighth

\item
For either sparse product prior, write the independent cell intensities as
\[
\mathbf Z^{\mathrm{sp}}
=(Z_x^{\mathrm{sp}})_{x\in[d]}
\sim\nu_\iota^{\otimes d},
\qquad
\mathbf Z^{\mathrm{sp}}\in[0,M]^d,
\]
and define
\[
H_{\mathrm{sp}}
=\frac1d\sum_{x\in[d]}\phi_\epsilon(Z_x^{\mathrm{sp}}),
\qquad
S_{\mathrm{sp}}
=1-\epsilon+\frac{\epsilon}{d}\sum_{x\in[d]}Z_x^{\mathrm{sp}},
\]
\[
\theta_{\mathrm{sp},\iota}
=\frac12+\frac12\int\phi_\epsilon\,d\nu_\iota,
\qquad
\delta_{\mathrm{sp}}
=|\theta_{\mathrm{sp},1}-\theta_{\mathrm{sp},0}|
=\frac12\left|\int\phi_\epsilon\,d(\nu_1-\nu_0)\right|.
\]
The sparse value identity in
\cref{obj:lem:sparse-likelihood} is
\[
\Psi(\mathbb P_{\mathbf Z^{\mathrm{sp}}}^{\mathrm{sp}})
=\frac12+\frac{H_{\mathrm{sp}}}{2S_{\mathrm{sp}}}.
\]
The mean-one constraint and \(0\le Z_x^{\mathrm{sp}}\le M\) give
\[
E (Z_x^{\mathrm{sp}})^2\le M,
\qquad
E S_{\mathrm{sp}}=1,
\qquad
E(S_{\mathrm{sp}}-1)^2\le\frac{\epsilon^2M}{d}.
\]
For \(0\le z\le M\), both \(z-1\) and zero are at most \(1+z\), so
\[
0\le(z-1)_+\le1+z,
\qquad
0\le1-\epsilon+\epsilon z,
\qquad
1+z>0.
\]
Multiplying the first inequality by the nonnegative factor
\((1-\epsilon+\epsilon z)/(1+z)\) gives
\[
0\le\phi_\epsilon(z)\le1-\epsilon+\epsilon z.
\]
Moreover, \(z^2\le Mz\), \((1-\epsilon)^2\le1\), and the nonnegativity of
\((1-\epsilon-\epsilon z)^2\) yields
\[
\begin{aligned}
\phi_\epsilon(z)^2
&\le(1-\epsilon+\epsilon z)^2\\
&\le2\{(1-\epsilon)^2+\epsilon^2z^2\}\\
&\le2(1+\epsilon^2Mz).
\end{aligned}
\]
Independence therefore yields
\[
E\{H_{\mathrm{sp}}-EH_{\mathrm{sp}}\}^2
\le\frac{2(1+\epsilon^2M)}d.
\]
Since \(S_{\mathrm{sp}}\ge1-\epsilon\ge1/2\) and
\(0\le H_{\mathrm{sp}}\le S_{\mathrm{sp}}\),
\[
\left|\frac{H_{\mathrm{sp}}}{S_{\mathrm{sp}}}
-H_{\mathrm{sp}}\right|
\le|1-S_{\mathrm{sp}}|.
\]
Using \((a+b)^2\le2(a^2+b^2)\) now gives
\[
E\left[
\left\{\Psi(\mathbb P_{\mathbf Z^{\mathrm{sp}}}^{\mathrm{sp}})
-\theta_{\mathrm{sp},\iota}\right\}^2
\right]
\le\frac{2+3\epsilon^2M}{2d}.
\]
Consequently, if
\[
64(2+3\epsilon^2M)\le d\delta_{\mathrm{sp}}^2,
\]
Markov's inequality gives
\[
\Pr_\iota\left(
\left|\Psi(\mathbb P_{\mathbf Z^{\mathrm{sp}}}^{\mathrm{sp}})
-\theta_{\mathrm{sp},\iota}\right|
>\frac{\delta_{\mathrm{sp}}}{4}
\right)\le\frac18.
\]

For completeness, define the two integrated squared risks of a count estimator \(\widehat W\) by
\[
\mathcal B_{\mathrm{sp},\iota}(\widehat W)
=
E_{\mathbf Z^{\mathrm{sp}}\sim\nu_\iota^{\otimes d}}
E_{\mathrm{counts}\mid\mathbf Z^{\mathrm{sp}}}
\left[
\{\widehat W-\Psi(\mathbb P_{\mathbf Z^{\mathrm{sp}}}^{\mathrm{sp}})\}^2
\right].
\]
If \(\theta_{\mathrm{sp},0}\le\theta_{\mathrm{sp},1}\), threshold
\(\widehat W\) at
\[
\frac{\theta_{\mathrm{sp},0}+\theta_{\mathrm{sp},1}}2.
\]
Otherwise interchange the prior indices and use the same midpoint. On an erroneous decision with the target within
\(\delta_{\mathrm{sp}}/4\) of its center, estimation error is at least
\(\delta_{\mathrm{sp}}/4\). The sum of the two testing errors is at least one minus their predictive total variation. Subtracting the two target-deviation probabilities thus proves
\[
\begin{aligned}
\max_{\iota=0,1}\mathcal B_{\mathrm{sp},\iota}(\widehat W)
&\ge
\frac12\left(\frac{\delta_{\mathrm{sp}}}{4}\right)^2
\left(1-\frac18-\frac18-\frac18\right)\\
&=\frac{5\delta_{\mathrm{sp}}^2}{256}.
\end{aligned}
\]

Use the completed packet pair, whose gap is at least
\(c_{\mathrm p}\sqrt M/K_{\mathrm{test}}\). Then
\[
\delta_{\mathrm{sp}}^2
\ge\frac{c_{\mathrm p}^2M}{4K_{\mathrm{test}}^2}.
\]
Thus the condition
\[
256(2+3\epsilon^2M)K_{\mathrm{test}}^2
\le c_{\mathrm p}^2Md
\]
supplies the target-deviation bound, and, with
\[
c_{\mathrm{sp}}=\frac{5c_{\mathrm p}^2}{1024},
\]
every measurable count estimator satisfies
\[
\max_{\iota=0,1}\mathcal B_{\mathrm{sp},\iota}(\widehat W)
\ge c_{\mathrm{sp}}\frac{M}{K_{\mathrm{test}}^2}.
\]
% lean: CausalSmith.Stat.OptvalueVanishingoverlapRate.sparse_packet_bayesRisk_lower

\item
We record the logarithmic degree and alphabet calibration explicitly. For real \(C_{\mathrm{deg}}\ge2\) and integer \(d\ge2\), define
\[
\mathcal K_{\mathrm{cal}}(C_{\mathrm{deg}},d)
=2\left\lceil\frac{C_{\mathrm{deg}}L_d}{2}\right\rceil.
\]
The ceiling inequalities and \(L_d\ge1\) give
\[
2\le\mathcal K_{\mathrm{cal}}(C_{\mathrm{deg}},d),
\qquad
C_{\mathrm{deg}}L_d
\le\mathcal K_{\mathrm{cal}}(C_{\mathrm{deg}},d)
\le(C_{\mathrm{deg}}+2)L_d.
\]
Choose
\[
C_{\mathrm{sp}}=\max\{C_{\mathrm{tv}},2\},
\qquad
K_{\mathrm{sp}}=\mathcal K_{\mathrm{cal}}(C_{\mathrm{sp}},d).
\]
Because \(L_d^2/d\to0\), a universal integer \(D_{\mathrm{sp}}\ge2\) can be chosen so that, for \(d\ge D_{\mathrm{sp}}\),
\[
\frac{L_d^2}{d}
<
\frac{c_{\mathrm p}^2}{448(C_{\mathrm{sp}}+2)^2},
\qquad
448K_{\mathrm{sp}}^2\le c_{\mathrm p}^2d.
\]
The asserted logarithmic limit follows, for example, by writing
\(d=e^{L_d-1}\) and using that an exponential dominates a fixed power. For \(M\ge2\) and \(0\le\epsilon\le1/2\),
\[
256(2+3\epsilon^2M)\le448M.
\]
Hence this single cutoff ensures
\[
256(2+3\epsilon^2M)K_{\mathrm{sp}}^2
\le c_{\mathrm p}^2Md
\]
uniformly in the required support and overlap ranges. The previous two steps therefore apply whenever
\[
2\le M\le K_{\mathrm{sp}}^2,
\qquad
M\le\frac{\eta_{\mathrm{sp}}dK_{\mathrm{sp}}}{Bn\epsilon}.
\]
% lean: CausalSmith.Stat.OptvalueVanishingoverlapRate.sparse_packet_bayesRisk_lower_largeAlphabet

\item
\textbf{Dense comparison.}\quad
The dense construction uses symmetric moment-matched priors on \([-1,1]\). We establish the absolute-moment gap used for them.

For \(K_{\mathrm{abs}}\ge1\), define
\[
E_{\mathrm{abs}}(K_{\mathrm{abs}})
=\inf_{\deg P_{\mathrm{abs}}\le K_{\mathrm{abs}}}
\sup_{-1\le z\le1}\bigl||z|-P_{\mathrm{abs}}(z)\bigr|.
\]
On the circle of period \(\pi\), define, for each integer \(N_{\mathrm{Fej}}\ge1\),
\[
F_{\mathrm{Fej},N_{\mathrm{Fej}}}(\omega)
=\frac1{N_{\mathrm{Fej}}}
\left|\sum_{\ell_{\mathrm{Fej}}=0}^{N_{\mathrm{Fej}}-1}
 e^{2i\ell_{\mathrm{Fej}}\omega}\right|^2,
\qquad \omega\in\mathbb R,
\]
\[
V_{\mathrm{VP},K_{\mathrm{abs}}}
=2F_{\mathrm{Fej},2K_{\mathrm{abs}}}
-F_{\mathrm{Fej},K_{\mathrm{abs}}}.
\]
For a continuous complex-valued \(\pi\)-periodic function \(f_{\mathrm{cusp}}\), define
\[
\Lambda_{\mathrm{cusp},K_{\mathrm{abs}}}(f_{\mathrm{cusp}})
=f_{\mathrm{cusp}}(0)-\int_0^\pi
V_{\mathrm{VP},K_{\mathrm{abs}}}(\omega)
 f_{\mathrm{cusp}}(\omega)\,\frac{d\omega}{\pi}.
\]
Expanding the squared finite sum gives
\[
F_{\mathrm{Fej},N_{\mathrm{Fej}}}(\omega)
=\frac1{N_{\mathrm{Fej}}}
\sum_{\ell_{\mathrm{Fej}},r_{\mathrm{Fej}}=0}^{N_{\mathrm{Fej}}-1}
 e^{2i(\ell_{\mathrm{Fej}}-r_{\mathrm{Fej}})\omega}.
\]
For each integer frequency \(j_{\mathrm{Four}}\), orthogonality retains exactly the pairs whose index difference is \(-j_{\mathrm{Four}}\). There are
\(N_{\mathrm{Fej}}-|j_{\mathrm{Four}}|\) such pairs when
\(|j_{\mathrm{Four}}|<N_{\mathrm{Fej}}\), and none otherwise. Thus the triangular coefficients are
\[
\begin{aligned}
\gamma_{\mathrm{Fej}}(N_{\mathrm{Fej}},j_{\mathrm{Four}})
&:=\int_0^\pi F_{\mathrm{Fej},N_{\mathrm{Fej}}}(\omega)
 e^{2ij_{\mathrm{Four}}\omega}\,\frac{d\omega}{\pi}\\
&=\begin{cases}
1-|j_{\mathrm{Four}}|/N_{\mathrm{Fej}},
&|j_{\mathrm{Four}}|<N_{\mathrm{Fej}},\\
0,&|j_{\mathrm{Four}}|\ge N_{\mathrm{Fej}}.
\end{cases}
\end{aligned}
\]
In particular, both Fej\'er kernels are nonnegative with integral one, so
\[
|\Lambda_{\mathrm{cusp},K_{\mathrm{abs}}}(f_{\mathrm{cusp}})|
\le4\|f_{\mathrm{cusp}}\|_\infty.
\]
Substituting the triangular coefficients into the cusp functional gives
\[
\begin{aligned}
\mu_{\mathrm{cusp}}(j_{\mathrm{Four}})
&:=\Lambda_{\mathrm{cusp},K_{\mathrm{abs}}}
 (\omega\mapsto e^{2ij_{\mathrm{Four}}\omega})\\
&=1-2\gamma_{\mathrm{Fej}}(2K_{\mathrm{abs}},j_{\mathrm{Four}})
 +\gamma_{\mathrm{Fej}}(K_{\mathrm{abs}},j_{\mathrm{Four}})\\
&=\begin{cases}
0,&|j_{\mathrm{Four}}|\le K_{\mathrm{abs}},\\
|j_{\mathrm{Four}}|/K_{\mathrm{abs}}-1,
&K_{\mathrm{abs}}<|j_{\mathrm{Four}}|<2K_{\mathrm{abs}},\\
1,&|j_{\mathrm{Four}}|\ge2K_{\mathrm{abs}}.
\end{cases}
\end{aligned}
\]
These multipliers are nonnegative. The identity
\[
\sin2\omega=\frac{e^{2i\omega}-e^{-2i\omega}}{2i}
\]
and the binomial expansion show that a polynomial of degree at most
\(K_{\mathrm{abs}}\) in \(\sin2\omega\) has frequencies
\(|j_{\mathrm{Four}}|\le K_{\mathrm{abs}}\). Hence
\[
\Lambda_{\mathrm{cusp},K_{\mathrm{abs}}}
 (\omega\mapsto P_{\mathrm{abs}}(\sin2\omega))=0
\qquad(\deg P_{\mathrm{abs}}\le K_{\mathrm{abs}}).
\]

To compute the absolute-sine coefficients, define
\[
S_{\mathrm{abs}}(\omega)=|\sin\omega|,
\qquad
c_{\mathrm{Four},\mathrm{abs}}(j_{\mathrm{Four}})
=\int_0^\pi S_{\mathrm{abs}}(\omega)e^{-2ij_{\mathrm{Four}}\omega}
 \,\frac{d\omega}{\pi},
\qquad j_{\mathrm{Four}}\in\mathbb Z.
\]
On \([0,\pi]\), \(S_{\mathrm{abs}}(\omega)=\sin\omega\). For every integer
\(j_{\mathrm{Four}}\), the denominator \(1-4j_{\mathrm{Four}}^2\) is nonzero, and the function
\[
\mathcal A_{\mathrm{Four},j_{\mathrm{Four}}}(\omega)
=e^{-2ij_{\mathrm{Four}}\omega}
\frac{-2ij_{\mathrm{Four}}\sin\omega-\cos\omega}
 {1-4j_{\mathrm{Four}}^2},
\qquad \omega\in[0,\pi],
\]
satisfies
\[
\mathcal A_{\mathrm{Four},j_{\mathrm{Four}}}'(\omega)
=e^{-2ij_{\mathrm{Four}}\omega}\sin\omega.
\]
Evaluating at the endpoints, using
\(e^{-2ij_{\mathrm{Four}}\pi}=1\), gives
\[
c_{\mathrm{Four},\mathrm{abs}}(j_{\mathrm{Four}})
=\frac{\mathcal A_{\mathrm{Four},j_{\mathrm{Four}}}(\pi)
 -\mathcal A_{\mathrm{Four},j_{\mathrm{Four}}}(0)}{\pi}
=-\frac2{\pi(4j_{\mathrm{Four}}^2-1)}.
\]
The zero coefficient is \(2/\pi\). For \(|j_{\mathrm{Four}}|\ge1\),
\[
|c_{\mathrm{Four},\mathrm{abs}}(j_{\mathrm{Four}})|
=\frac2{\pi(4j_{\mathrm{Four}}^2-1)}
\le\frac1{j_{\mathrm{Four}}^2},
\qquad
\sum_{j_{\mathrm{Four}}\in\mathbb Z}
 |c_{\mathrm{Four},\mathrm{abs}}(j_{\mathrm{Four}})|<\infty.
\]
Thus the Fourier series converges absolutely and uniformly; its coefficients identify its sum with the continuous function \(S_{\mathrm{abs}}\). Pairing opposite frequencies and replacing \(\omega\) by \(2\omega\) yields
\[
|\sin2\omega|
=\frac2\pi-\frac4\pi
 \sum_{\ell=1}^\infty\frac{\cos(4\ell\omega)}{4\ell^2-1}.
\]
Continuity of the cusp functional permits termwise application, giving
\[
-\Lambda_{\mathrm{cusp},K_{\mathrm{abs}}}(|\sin2\cdot|)
=\sum_{\ell=1}^\infty
 \frac{4\mu_{\mathrm{cusp}}(2\ell)}{\pi(4\ell^2-1)}.
\]
Every summand is nonnegative. For
\(K_{\mathrm{abs}}\le\ell<2K_{\mathrm{abs}}\), the multiplier equals one, and \(\pi\le4\) gives
\[
\frac4{\pi(4\ell^2-1)}\ge\frac1{16K_{\mathrm{abs}}^2}.
\] Keeping the
\(K_{\mathrm{abs}}\) terms with
\(K_{\mathrm{abs}}\le\ell<2K_{\mathrm{abs}}\) gives
\[
-\Lambda_{\mathrm{cusp},K_{\mathrm{abs}}}(|\sin2\cdot|)
\ge
\sum_{\ell=K_{\mathrm{abs}}}^{2K_{\mathrm{abs}}-1}
\frac4{\pi(4\ell^2-1)}
\ge\frac1{16K_{\mathrm{abs}}}.
\]
It follows that
\[
\frac1{16K_{\mathrm{abs}}}
\le
4\sup_{-1\le z\le1}\bigl||z|-P_{\mathrm{abs}}(z)\bigr|,
\]
and therefore
\[
E_{\mathrm{abs}}(K_{\mathrm{abs}})
\ge\frac1{100K_{\mathrm{abs}}}.
\]

To construct the priors, let
\[
\mathcal V_{\mathrm{abs}}
=\{P_{\mathrm{abs}}|_{[-1,1]}:\deg P_{\mathrm{abs}}\le K_{\mathrm{abs}}\}.
\]
The norm of the class of \(|z|\) in
\(C([-1,1],\mathbb R)/\mathcal V_{\mathrm{abs}}\) is
\(E_{\mathrm{abs}}(K_{\mathrm{abs}})\). A norming functional on that quotient, pulled back to \(C([-1,1],\mathbb R)\), supplies a functional
\(\Lambda_{\mathrm{abs}}\) satisfying
\[
\|\Lambda_{\mathrm{abs}}\|\le1,
\qquad
\Lambda_{\mathrm{abs}}(z^\ell)=0
\quad(0\le\ell\le K_{\mathrm{abs}}),
\qquad
\Lambda_{\mathrm{abs}}(|z|)=E_{\mathrm{abs}}(K_{\mathrm{abs}}).
\]
Its positive half-split can be obtained by positive extension. In detail, let \(1\) denote the constant function with value one and consider the cone
\[
\mathcal C_{\mathrm{split}}
=\{(f_++f_-,-\Lambda_{\mathrm{abs}}(f_-)):
 f_+,f_-\in C([-1,1],\mathbb R),\ f_+,f_-\ge0\}.
\]
On the subspace of pairs \((a_{\mathrm{cone}}\,1,b_{\mathrm{cone}})\), prescribe the functional
\[
(a_{\mathrm{cone}}\,1,b_{\mathrm{cone}})
\longmapsto a_{\mathrm{cone}}/2+b_{\mathrm{cone}},
\qquad a_{\mathrm{cone}},b_{\mathrm{cone}}\in\mathbb R.
\]
If such a pair belongs to the cone, then
\[
f_++f_-=a_{\mathrm{cone}}\,1,
\qquad
b_{\mathrm{cone}}=-\Lambda_{\mathrm{abs}}(f_-),
\qquad a_{\mathrm{cone}}\ge0.
\]
Consequently,
\[
\|2f_--a_{\mathrm{cone}}\,1\|_\infty\le a_{\mathrm{cone}},
\qquad
\Lambda_{\mathrm{abs}}(f_-)
=\tfrac12\Lambda_{\mathrm{abs}}(2f_--a_{\mathrm{cone}}\,1)
\le a_{\mathrm{cone}}/2,
\]
where the last bound uses \(\|\Lambda_{\mathrm{abs}}\|\le1\) and
\(\Lambda_{\mathrm{abs}}(1)=0\). Thus the prescribed functional is nonnegative on its intersection with the cone.

The subspace also satisfies the required majorization condition. For every
\(f_{\mathrm{split}}\in C([-1,1],\mathbb R)\) and
\(t_{\mathrm{split}}\in\mathbb R\),
\[
(\|f_{\mathrm{split}}\|_\infty\,1,-t_{\mathrm{split}})
 +(f_{\mathrm{split}},t_{\mathrm{split}})
=(\|f_{\mathrm{split}}\|_\infty\,1+f_{\mathrm{split}},0)
\in\mathcal C_{\mathrm{split}}.
\]
Indeed, \(\|f_{\mathrm{split}}\|_\infty\,1+f_{\mathrm{split}}\ge0\), so this is the cone element obtained with that function as \(f_+\) and zero as \(f_-\). Positive extension therefore supplies a linear functional
\[
G_{\mathrm{split}}:C([-1,1],\mathbb R)\times\mathbb R\longrightarrow\mathbb R
\]
that is nonnegative on \(\mathcal C_{\mathrm{split}}\) and satisfies
\[
G_{\mathrm{split}}(a_{\mathrm{cone}}\,1,b_{\mathrm{cone}})
=a_{\mathrm{cone}}/2+b_{\mathrm{cone}}.
\]
Define, for every \(f_{\mathrm{split}}\in C([-1,1],\mathbb R)\),
\[
\Lambda_+(f_{\mathrm{split}})=G_{\mathrm{split}}(f_{\mathrm{split}},0),
\qquad
\Lambda_-(f_{\mathrm{split}})
=\Lambda_+(f_{\mathrm{split}})-\Lambda_{\mathrm{abs}}(f_{\mathrm{split}}).
\]
For \(f_{\mathrm{split}}\ge0\), both
\((f_{\mathrm{split}},0)\) and
\((f_{\mathrm{split}},-\Lambda_{\mathrm{abs}}(f_{\mathrm{split}}))\) belong to the cone. Since
\(G_{\mathrm{split}}(0,t_{\mathrm{split}})=t_{\mathrm{split}}\), this gives
\[
\Lambda_+(f_{\mathrm{split}})\ge0,
\qquad
\Lambda_-(f_{\mathrm{split}})
=G_{\mathrm{split}}(f_{\mathrm{split}},-\Lambda_{\mathrm{abs}}(f_{\mathrm{split}}))\ge0.
\]
Their masses and difference are
\[
\Lambda_+(1)=G_{\mathrm{split}}(1,0)=\frac12,
\qquad
\Lambda_-(1)=\Lambda_+(1)-\Lambda_{\mathrm{abs}}(1)=\frac12,
\qquad
\Lambda_+-\Lambda_-=\Lambda_{\mathrm{abs}}.
\]
Positivity bounds each functional on a continuous function by its mass times the uniform norm, so both have representing finite measures.
Their representing measures have mass \(1/2\). Averaging each measure with its reflection under \(z\mapsto-z\) preserves the absolute-value pairing and the annihilation of all powers through degree \(K_{\mathrm{abs}}\). Doubling these symmetric measures gives symmetric probability measures
\(\eta_0,\eta_1\) on \([-1,1]\) satisfying
\[
\int z^\ell\,d\eta_0=\int z^\ell\,d\eta_1
\quad(0\le\ell\le K_{\mathrm{abs}}),
\]
\[
\int|z|\,d\eta_1-\int|z|\,d\eta_0
=2E_{\mathrm{abs}}(K_{\mathrm{abs}})
\ge\frac1{50K_{\mathrm{abs}}}.
\]
% lean: CausalSmith.Stat.OptvalueVanishingoverlapRate.dense_symmetric_priors_absGap_lower

\item
For \(0<\mathfrak a\le1/2\), define the independent prior draws by
\[
\mathbf T^{\mathrm{de}}
=(T_x^{\mathrm{de}})_{x\in[d]}
\sim\eta_\iota^{\otimes d},
\qquad
\mathbf T^{\mathrm{de}}\in[-1,1]^d,
\]
and define the dense observed law by
\[
p_x=\frac1d,
\qquad
\pi_x=\epsilon,
\qquad
\mu_{0x}=\frac12,
\qquad
\mu_{1x}=\frac12+\mathfrak a T_x^{\mathrm{de}}.
\]
Equivalently, its four atom probabilities are
\[
q_{00,x}=q_{01,x}=\frac{1-\epsilon}{2d},
\qquad
q_{10,x}=\frac{\epsilon}{d}\left(\frac12-\mathfrak a T_x^{\mathrm{de}}\right),
\qquad
q_{11,x}=\frac{\epsilon}{d}\left(\frac12+\mathfrak a T_x^{\mathrm{de}}\right).
\]
These are legal observed laws, including the endpoints \(T_x^{\mathrm{de}}=\pm1\).
Their target and prior centers are
\[
\Psi_{\mathrm{de}}(\mathbf T^{\mathrm{de}})
=\frac12+\frac{\mathfrak a}{2d}
 \sum_{x\in[d]}(|T_x^{\mathrm{de}}|+T_x^{\mathrm{de}}),
\]
\[
\theta_{\mathrm{de},\iota}
=\frac12+\frac{\mathfrak a}{2}\int|z|\,d\eta_\iota,
\qquad
\delta_{\mathrm{de}}
=\frac{\mathfrak a}{2}
 \left(\int|z|\,d\eta_1-\int|z|\,d\eta_0\right)
\ge\frac{\mathfrak a}{100K_{\mathrm{abs}}}.
\]
Symmetry gives \(\int z\,d\eta_\iota=0\). Since
\((|z|+z)/2=z_+\in[0,1]\), independence gives
\[
E_\iota\{\Psi_{\mathrm{de}}-\theta_{\mathrm{de},\iota}\}^2
\le\frac{\mathfrak a^2}{d}.
\]

In the Poisson experiment, define the total treated intensity per cell by
\[
\lambda_{\mathrm{tr}}=\frac{Bn\epsilon}{d}.
\]
The treated counts have means
\[
\lambda_{\mathrm{tr}}\left(\frac12-\mathfrak a z\right),
\qquad
\lambda_{\mathrm{tr}}\left(\frac12+\mathfrak a z\right).
\]
For nonnegative integer treated failure and success counts, respectively denoted by \(N_{\mathrm{de},0}\) and \(N_{\mathrm{de},1}\), define the likelihood ratio relative to \(z=0\) by
\[
\begin{aligned}
L_{\mathrm{de},z}(N_{\mathrm{de},0},N_{\mathrm{de},1})
&=e^{\lambda_{\mathrm{tr}}\mathfrak a z}
 (1-2\mathfrak a z)^{N_{\mathrm{de},0}}
 e^{-\lambda_{\mathrm{tr}}\mathfrak a z}
 (1+2\mathfrak a z)^{N_{\mathrm{de},1}}\\
&=(1-2\mathfrak a z)^{N_{\mathrm{de},0}}
  (1+2\mathfrak a z)^{N_{\mathrm{de},1}},
\qquad z\in[-1,1].
\end{aligned}
\]
At a zero alternative mean, use \(0^0=1\); a zero base raised to a positive count is zero. The exponential factors cancel because the sum of the two treated means is \(\lambda_{\mathrm{tr}}\). Multiplication of these likelihood factors and their expectations gives the Gram identity
\[
E_0[L_{\mathrm{de},z}L_{\mathrm{de},w}]
=\exp\{4\lambda_{\mathrm{tr}}\mathfrak a^2zw\},
\qquad z,w\in[-1,1].
\]
Thus, defining
\[
\tau_{\mathrm{de}}=4\lambda_{\mathrm{tr}}\mathfrak a^2,
\qquad
\eta_{\mathrm{de}}=\frac{\beta_{\mathrm{tail}}}{4},
\]
the same moment-matching and Cauchy--Schwarz argument gives predictive total variation at most
\[
d\sqrt{\mathcal T_{K_{\mathrm{abs}}}(\tau_{\mathrm{de}})}
\le\frac18
\]
whenever
\[
K_{\mathrm{abs}}\ge C_{\mathrm{tv}}L_d,
\qquad
\lambda_{\mathrm{tr}}\mathfrak a^2
\le\eta_{\mathrm{de}}K_{\mathrm{abs}}.
\]
The two untreated counts have common means
\(Bn(1-\epsilon)/(2d)\), independent of the prior parameter. Adding them preserves this total-variation bound, so it applies to all four observed counts.

If
\[
128\mathfrak a^2\le d\delta_{\mathrm{de}}^2,
\]
the target-deviation probability at radius
\(\delta_{\mathrm{de}}/4\) is at most \(1/8\). Applying the midpoint argument above to the two dense centers gives, for every measurable estimator of the full counts,
\[
\max_{\iota=0,1}\mathcal B_{\mathrm{de},\iota}(\widehat W)
\ge\frac{5\delta_{\mathrm{de}}^2}{256}
\ge\frac5{2560000}\frac{\mathfrak a^2}{K_{\mathrm{abs}}^2},
\]
where
\[
\mathcal B_{\mathrm{de},\iota}(\widehat W)
=
E_{\mathbf T^{\mathrm{de}}\sim\eta_\iota^{\otimes d}}
E_{\mathrm{counts}\mid\mathbf T^{\mathrm{de}}}
[\{\widehat W-\Psi_{\mathrm{de}}(\mathbf T^{\mathrm{de}})\}^2].
\]
Choose
\[
C_{\mathrm{de}}=\max\{C_{\mathrm{tv}},2\},
\qquad
K_{\mathrm{de}}=\mathcal K_{\mathrm{cal}}(C_{\mathrm{de}},d).
\]
The same logarithmic cutoff construction, applied with packet comparison constant \(1/100\), supplies a universal
\(D_{\mathrm{de}}\ge2\) such that
\[
1280000K_{\mathrm{de}}^2\le d
\quad(d\ge D_{\mathrm{de}}).
\]
Since
\(\delta_{\mathrm{de}}^2\ge\mathfrak a^2/(10000K_{\mathrm{de}}^2)\),
this implies the required target-deviation condition. Reordering the four counts into the observed atom histogram is a bijection, so both integrated-risk bounds hold on that histogram space.
% lean: CausalSmith.Stat.OptvalueVanishingoverlapRate.dense_full_observed_logDegree_bayesRisk_lower

\item
We transfer these bounds to fixed samples. Let \(\widehat V\) be any measurable estimator based on \(n\) observations. For a histogram
\[
\boldsymbol h_{\mathrm{hist}}:\mathcal J_d\to\mathbb N,
\qquad
n_{\mathrm{hist}}=\sum_{o\in\mathcal J_d}\boldsymbol h_{\mathrm{hist}}(o),
\]
define
\[
\mathcal S(\boldsymbol h_{\mathrm{hist}})
=\{\boldsymbol s\in\mathcal J_d^{n_{\mathrm{hist}}}:
\boldsymbol s\text{ has histogram }\boldsymbol h_{\mathrm{hist}}\},
\qquad
\boldsymbol o^\circ=((0,0,0),\ldots,(0,0,0))\in\mathcal J_d^n,
\]
\[
\boldsymbol o_n(\boldsymbol s)=
\begin{cases}
(s_1,\ldots,s_n),&n_{\mathrm{hist}}\ge n,\\
\boldsymbol o^\circ,&n_{\mathrm{hist}}<n,
\end{cases}
\]
\[
\widehat W_{\widehat V}(\boldsymbol h_{\mathrm{hist}})
=\frac1{|\mathcal S(\boldsymbol h_{\mathrm{hist}})|}
 \sum_{\boldsymbol s\in\mathcal S(\boldsymbol h_{\mathrm{hist}})}
 \Pi_{[0,1]}\{\widehat V(\boldsymbol o_n(\boldsymbol s))\}.
\]
The set \(\mathcal S(\boldsymbol h_{\mathrm{hist}})\) is nonempty, including the zero histogram, for which it contains the empty string.

For independent Poisson atom counts with total mean
\(\lambda_{\mathrm{tot}}\), conditional on their total, the labels are independent with the normalized observed law; conditional on the histogram, their ordering is uniform on \(\mathcal S(h)\). Jensen's inequality, followed by projection onto \([0,1]\), therefore gives
\[
E\{\widehat W_{\widehat V}-\Psi\}^2
\le E\{\widehat V-\Psi\}^2
+\Pr\{\operatorname{Pois}(\lambda_{\mathrm{tot}})<n\}.
\]
For the sparse experiment,
\[
\lambda_{\mathrm{tot}}=BnS_{\mathrm{sp}}\ge Bn/2;
\]
for the dense experiment, \(\lambda_{\mathrm{tot}}=Bn\).

For any \(\lambda_{\mathrm{tot}}\ge Bn/2>n\), put
\[
\gamma_{\mathrm{sh}}
=\frac{\lambda_{\mathrm{tot}}-n}{\lambda_{\mathrm{tot}}}\in(0,1).
\]
Exponential Markov inequality and
\(e^{-\gamma_{\mathrm{sh}}}
\le1-\gamma_{\mathrm{sh}}+\gamma_{\mathrm{sh}}^2/2\) give
\[
\Pr\{\operatorname{Pois}(\lambda_{\mathrm{tot}})<n\}
\le
\exp\left\{-\frac{(\lambda_{\mathrm{tot}}-n)^2}
                       {2\lambda_{\mathrm{tot}}}\right\}
\le
\exp\left\{-\frac{n(B-2)^2}{4B}\right\}.
\]
The last inequality uses that
\((\lambda_{\mathrm{tot}}-n)^2/(2\lambda_{\mathrm{tot}})\)
is increasing for \(\lambda_{\mathrm{tot}}>n\).

Define the shortage penalty by
\[
p_{\mathrm{sh}}(n,B)
=\exp\left\{-\frac{n(B-2)^2}{4B}\right\}.
\]
Integrating the reconstruction bound under either supported prior and then taking the fixed-sample estimator infimum proves
\[
c_{\mathrm{sp}}\frac M{K_{\mathrm{sp}}^2}
-p_{\mathrm{sh}}(n,B)
\le\mathfrak R_{n,d,\epsilon}
\]
under the sparse conditions, and
\[
\frac5{2560000}\frac{\mathfrak a^2}{K_{\mathrm{de}}^2}
-p_{\mathrm{sh}}(n,B)
\le\mathfrak R_{n,d,\epsilon}
\]
under the dense conditions.

For the sparse bound choose
\[
M_{\mathrm{sp}}
=\min\left\{
K_{\mathrm{sp}}^2,
\frac{\eta_{\mathrm{sp}}dK_{\mathrm{sp}}}{Bn\epsilon}
\right\}.
\]
Whenever the second entry is at least two, this is an admissible support endpoint, and
\[
c_{\mathrm{sp}}
\min\left\{1,\frac{\eta_{\mathrm{sp}}d}
                    {Bn\epsilon K_{\mathrm{sp}}}\right\}
-p_{\mathrm{sh}}(n,B)
\le\mathfrak R_{n,d,\epsilon}.
\]
For the dense bound choose
\[
\mathfrak a_{\mathrm{de}}^2
=\min\left\{\frac14,
\frac{\eta_{\mathrm{de}}dK_{\mathrm{de}}}{Bn\epsilon}\right\}.
\]
Then \(0<\mathfrak a_{\mathrm{de}}\le1/2\) and the likelihood condition holds. Dividing by \(K_{\mathrm{de}}^2\) gives
\[
\frac5{2560000}
\min\left\{
\frac1{4K_{\mathrm{de}}^2},
\frac{\eta_{\mathrm{de}}d}{Bn\epsilon K_{\mathrm{de}}}
\right\}
-p_{\mathrm{sh}}(n,B)
\le\mathfrak R_{n,d,\epsilon}.
\]
% lean: CausalSmith.Stat.OptvalueVanishingoverlapRate.sparse_fixedSample_calibrated_lower

\item
\textbf{Fixed-sample transfer.}\quad
We join these two ranges while retaining their degree comparison. Define
\[
D_{\mathrm{lg}}=\max\{D_{\mathrm{sp}},D_{\mathrm{de}}\},
\qquad
q_{\mathrm{deg}}=\frac{C_{\mathrm{de}}+2}{C_{\mathrm{sp}}},
\]
\[
b_{\mathrm{splice}}
=\min\left\{
\frac{\eta_{\mathrm{sp}}}{8q_{\mathrm{deg}}^2},
\frac{\eta_{\mathrm{de}}}{q_{\mathrm{deg}}}
\right\},
\qquad
c_{\mathrm{splice}}
=\min\left\{
c_{\mathrm{sp}}\min\{1,\eta_{\mathrm{sp}}\},
\frac5{2560000}b_{\mathrm{splice}}
\right\}>0,
\]
\[
x_{\mathrm{splice}}=\frac d{Bn\epsilon}.
\]
The degree bounds imply
\[
K_{\mathrm{de}}\le q_{\mathrm{deg}}K_{\mathrm{sp}}.
\]

If \(\eta_{\mathrm{sp}}x_{\mathrm{splice}}K_{\mathrm{sp}}\ge2\), use the sparse bound. The elementary inequality
\[
\min\{1,\eta_{\mathrm{sp}}\}
\min\left\{1,\frac{x_{\mathrm{splice}}}{K_{\mathrm{sp}}}\right\}
\le
\min\left\{1,
\frac{\eta_{\mathrm{sp}}x_{\mathrm{splice}}}{K_{\mathrm{sp}}}\right\}
\]
shows that its leading term is at least
\(c_{\mathrm{splice}}\min\{1,x_{\mathrm{splice}}/K_{\mathrm{sp}}\}\).

Otherwise,
\(\eta_{\mathrm{sp}}x_{\mathrm{splice}}K_{\mathrm{sp}}\le2\).
The degree comparison and the definition of \(b_{\mathrm{splice}}\) give
\[
b_{\mathrm{splice}}
\frac{x_{\mathrm{splice}}}{K_{\mathrm{sp}}}
\le\frac1{4K_{\mathrm{de}}^2},
\qquad
b_{\mathrm{splice}}
\frac{x_{\mathrm{splice}}}{K_{\mathrm{sp}}}
\le\frac{\eta_{\mathrm{de}}x_{\mathrm{splice}}}{K_{\mathrm{de}}}.
\]
For the first inequality, multiply by \(4K_{\mathrm{de}}^2\) and use
\[
\frac{\eta_{\mathrm{sp}}}{8q_{\mathrm{deg}}^2}
\frac{x_{\mathrm{splice}}}{K_{\mathrm{sp}}}
\,4K_{\mathrm{de}}^2
\le\frac{\eta_{\mathrm{sp}}x_{\mathrm{splice}}K_{\mathrm{sp}}}{2}
\le1.
\]
For the second, use
\(K_{\mathrm{de}}\le q_{\mathrm{deg}}K_{\mathrm{sp}}\).
The dense bound therefore supplies the same leading term. In both cases,
\[
c_{\mathrm{splice}}
\min\left\{1,\frac d{Bn\epsilon K_{\mathrm{sp}}}\right\}
-p_{\mathrm{sh}}(n,B)
\le\mathfrak R_{n,d,\epsilon}
\qquad(d\ge D_{\mathrm{lg}}).
\]

Set
\[
c_{\mathrm{sh}}=\frac{c_{\mathrm{splice}}}{C_{\mathrm{sp}}+2}.
\]
Since \(K_{\mathrm{sp}}\le(C_{\mathrm{sp}}+2)L_d\) and
\(1/\{B(C_{\mathrm{sp}}+2)\}\le1\),
\[
\frac1{B(C_{\mathrm{sp}}+2)}\rho_{n,d,\epsilon}
\le
\min\left\{1,\frac d{Bn\epsilon K_{\mathrm{sp}}}\right\}.
\]
Thus
\[
\frac{c_{\mathrm{sh}}}{B}\rho_{n,d,\epsilon}
-p_{\mathrm{sh}}(n,B)
\le\mathfrak R_{n,d,\epsilon}
\qquad(d\ge D_{\mathrm{lg}},\ B>2).
\]
% lean: CausalSmith.Stat.OptvalueVanishingoverlapRate.largeAlphabet_lower_with_shortage

\item
Choose once and for all
\[
B_{\mathrm{lg}}=\max\{4,2048/c_{\mathrm{sh}}\},
\qquad
c_{\mathrm{lg}}=\frac{c_{\mathrm{sh}}}{2B_{\mathrm{lg}}}>0.
\]
For \(B\ge4\), define
\[
a_{\mathrm{sh}}=\frac{(B-2)^2}{8B}\ge\frac B{32}.
\]
The inequality \(e^{-v}\le1/v\) for \(v>0\) gives
\[
e^{-a_{\mathrm{sh}}n}\le\frac{32}{Bn},
\qquad
e^{-a_{\mathrm{sh}}n}\le e^{-a_{\mathrm{sh}}}\le\frac{32}{B}.
\]
Multiplying these bounds proves
\[
p_{\mathrm{sh}}(n,B)
=e^{-a_{\mathrm{sh}}n}e^{-a_{\mathrm{sh}}n}
\le\frac{1024}{B^2n}.
\]
Also \(\log d\le d-1\), so \(L_d\le d\). Since \(n\ge1\) and
\(\epsilon\le1/2\le1\),
\[
\frac1n\le1,
\qquad
\frac1n\le\frac d{n\epsilon L_d},
\qquad
\frac1n\le\rho_{n,d,\epsilon}.
\]
Because \(B_{\mathrm{lg}}c_{\mathrm{sh}}\ge2048\),
\[
p_{\mathrm{sh}}(n,B_{\mathrm{lg}})
\le\frac{1024}{B_{\mathrm{lg}}^2n}
\le\frac{c_{\mathrm{sh}}}{2B_{\mathrm{lg}}}\frac1n
\le c_{\mathrm{lg}}\rho_{n,d,\epsilon}.
\]
The preceding lower bound has leading coefficient
\(c_{\mathrm{sh}}/B_{\mathrm{lg}}=2c_{\mathrm{lg}}\). Subtracting the shortage penalty therefore gives
\[
c_{\mathrm{lg}}\rho_{n,d,\epsilon}
\le\mathfrak R_{n,d,\epsilon}
\qquad(d\ge D_{\mathrm{lg}}).
\]
% lean: CausalSmith.Stat.OptvalueVanishingoverlapRate.largeAlphabet_minimaxRisk_lower

\item
For the bounded-alphabet comparison, define
\[
r_{\mathrm{par}}=\min\left\{1,\frac1{n\epsilon}\right\},
\qquad
a_{\mathrm{par}}=\sqrt{\frac{r_{\mathrm{par}}}{64}}.
\]
Then
\[
0<a_{\mathrm{par}}\le\frac12,
\qquad
n\epsilon a_{\mathrm{par}}^2\le\frac1{16}.
\]
Use two dense observed laws \(\mathbb P_{\mathrm{par},1}\) and
\(\mathbb P_{\mathrm{par},0}\) with uniform covariate probabilities, propensity
\(\epsilon\), control mean \(1/2\), and treated means respectively
\(1/2+a_{\mathrm{par}}\) and \(1/2\). Define their scalar targets by
\[
\theta_{\mathrm{par},1}
=\Psi(\mathbb P_{\mathrm{par},1})=\frac12+a_{\mathrm{par}},
\qquad
\theta_{\mathrm{par},0}
=\Psi(\mathbb P_{\mathrm{par},0})=\frac12.
\]
For every \(x\in[d]\), the atom probabilities of these laws are
\[
\begin{aligned}
\mathbb P_{\mathrm{par},1}(x,0,0)
&=\mathbb P_{\mathrm{par},1}(x,0,1)
=\mathbb P_{\mathrm{par},0}(x,0,0)
=\mathbb P_{\mathrm{par},0}(x,0,1)
=\frac{1-\epsilon}{2d},\\
\mathbb P_{\mathrm{par},0}(x,1,0)
&=\mathbb P_{\mathrm{par},0}(x,1,1)=\frac{\epsilon}{2d},\\
\mathbb P_{\mathrm{par},1}(x,1,0)
&=\frac{\epsilon}{d}\left(\frac12-a_{\mathrm{par}}\right),
\qquad
\mathbb P_{\mathrm{par},1}(x,1,1)
=\frac{\epsilon}{d}\left(\frac12+a_{\mathrm{par}}\right).
\end{aligned}
\]
All reference atom probabilities are positive. Each untreated atom has zero probability difference and hence zero chi-square contribution. Each treated atom contributes
\[
\frac{(\epsilon a_{\mathrm{par}}/d)^2}{\epsilon/(2d)}
=\frac{2\epsilon a_{\mathrm{par}}^2}{d}.
\]
Summing over both treated outcomes and all \(d\) cells gives the finite-atom calculation
\[
\begin{aligned}
\chi^2(\mathbb P_{\mathrm{par},1},\mathbb P_{\mathrm{par},0})
&=\sum_{x\in[d]}\sum_{a,y\in\{0,1\}}
\frac{\{\mathbb P_{\mathrm{par},1}(x,a,y)
-\mathbb P_{\mathrm{par},0}(x,a,y)\}^2}
 {\mathbb P_{\mathrm{par},0}(x,a,y)}\\
&=d\cdot2\cdot\frac{2\epsilon a_{\mathrm{par}}^2}{d}
=4\epsilon a_{\mathrm{par}}^2.
\end{aligned}
\]
Define the two \(n\)-observation laws by
\[
\mathbb Q_{\mathrm{par},\iota}
=\mathbb P_{\mathrm{par},\iota}^{\otimes n},
\qquad \iota\in\{0,1\}.
\]
Product likelihoods give
\[
1+\chi^2(\mathbb Q_{\mathrm{par},1},\mathbb Q_{\mathrm{par},0})
=(1+4\epsilon a_{\mathrm{par}}^2)^n
\le e^{4n\epsilon a_{\mathrm{par}}^2}
\le e^{1/4}\le\frac43.
\]
In particular their chi-square divergence is at most one, so their total variation is at most \(1/2\). Thresholding an arbitrary estimator at
\(1/2+a_{\mathrm{par}}/2\) shows that one of the two laws has probability at least \(1/4\) of absolute error at least \(a_{\mathrm{par}}/2\). Hence
\[
\max_{\iota=0,1}
E_{\mathbb Q_{\mathrm{par},\iota}}
\{\widehat V-\theta_{\mathrm{par},\iota}\}^2
\ge\frac{a_{\mathrm{par}}^2}{16}
=\frac1{1024}r_{\mathrm{par}}.
\]
Taking the estimator infimum gives
\[
\frac1{1024}\min\left\{1,\frac1{n\epsilon}\right\}
\le\mathfrak R_{n,d,\epsilon}
\qquad(d\ge2).
\]

Now define
\[
c_{\mathrm{all}}
=\min\left\{c_{\mathrm{lg}},\frac1{1024D_{\mathrm{lg}}}\right\}>0.
\]
For \(d\ge D_{\mathrm{lg}}\), the large-alphabet bound applies.
For \(2\le d<D_{\mathrm{lg}}\), using \(L_d\ge1\) gives
\[
\rho_{n,d,\epsilon}
\le
D_{\mathrm{lg}}\min\left\{1,\frac1{n\epsilon}\right\}.
\]
Indeed, \(\rho_{n,d,\epsilon}\le1\le D_{\mathrm{lg}}\), and
\[
\rho_{n,d,\epsilon}
\le\frac d{n\epsilon L_d}
\le\frac{D_{\mathrm{lg}}}{n\epsilon}.
\]
Thus in this second case,
\[
c_{\mathrm{all}}\rho_{n,d,\epsilon}
\le
\frac1{1024}\min\left\{1,\frac1{n\epsilon}\right\}
\le\mathfrak R_{n,d,\epsilon}.
\]
Together the two cases establish
\[
c_{\mathrm{all}}\rho_{n,d,\epsilon}
\le\mathfrak R_{n,d,\epsilon}
\qquad(n\ge1,\ d\ge2,\ 0<\epsilon\le1/2).
\]
% lean: CausalSmith.Stat.OptvalueVanishingoverlapRate.allAlphabet_minimaxRisk_lower

\item
\textbf{Final consequences.}\quad
By \cref{obj:thm:observable-upper}, there are universal calibrations
\(H_0>0,\kappa>0,D_0\ge2\) and a universal \(C_{\mathrm u}>0\) for the measurable estimator in
\cref{obj:def:armwise-estimator}. Define its worst-case observed risk by
\[
\mathcal R^{\mathrm{AF}}_{n,d,\epsilon}
=
\sup_{\mathbb P\in\mathcal D_{d,\epsilon}^{\mathrm{obs}}}
E_{\mathbb P}
\left[
\{\widehat V_{n,d,\epsilon}^{\mathrm{AF}}-\Psi(\mathbb P)\}^2
\right].
\]
Apply the cited bound to the product sampling law
\(\mathbb P^{\otimes n}\), which satisfies
\cref{obj:ass:iid-sampling}. This gives
\[
\mathcal R^{\mathrm{AF}}_{n,d,\epsilon}
\le C_{\mathrm u}\rho_{n,d,\epsilon}.
\]
Set
\[
c_\star=\min\{c_{\mathrm{all}},C_{\mathrm u}\},
\qquad
C_\star=C_{\mathrm u}.
\]
Then \(0<c_\star\le C_\star<\infty\). Nonnegativity of
\(\rho_{n,d,\epsilon}\) permits reducing the lower constant, and the estimator infimum in
\cref{obj:def:minimax-risk} is bounded by the risk of this particular estimator. Consequently,
\[
c_\star\rho_{n,d,\epsilon}
\le\mathfrak R_{n,d,\epsilon}
\le\mathcal R^{\mathrm{AF}}_{n,d,\epsilon}
\le C_\star\rho_{n,d,\epsilon}.
\]

For each measurable \(\widehat V\), identification and completion in
\cref{obj:prop:identification} give the equality
\[
\begin{aligned}
&\sup_{P\in\mathcal P_{d,\epsilon}}
E_{\operatorname{Obs}(P)}
[\{\widehat V-V^\star(P)\}^2]\\
&\hspace{2cm}=
\sup_{\mathbb P\in\mathcal D_{d,\epsilon}^{\mathrm{obs}}}
E_{\mathbb P}
[\{\widehat V-\Psi(\mathbb P)\}^2].
\end{aligned}
\]
For one inequality use
\(V^\star(P)=\Psi\{\operatorname{Obs}(P)\}\) and the observed marginal's membership in the observed class. For the other, complete each observed law to a causal law with that same marginal and value. Taking the estimator infimum preserves equality, so the observed bounds also prove the two stated causal minimax bounds.
% lean: CausalSmith.Stat.OptvalueVanishingoverlapRate.weighted_separation_and_frontier

\item
Finally, let the sequences in the statement be given. Define
\[
r_k=\frac{d_k}{n_k\epsilon_k\log(ed_k)},
\qquad
\rho_k=\min\{1,r_k\}.
\]
For the finitely many possible indices with \(n_k=0\), use the totalized division convention \(r_k=0\). Since \(n_k\to\infty\), one has \(n_k\ge1\) eventually. For all \(k\), \(r_k,\rho_k\ge0\), and
\[
\rho_k\longrightarrow0
\quad\Longleftrightarrow\quad
r_k\longrightarrow0.
\]
The reverse implication follows from \(0\le\rho_k\le r_k\).
For the forward implication, eventually \(\rho_k<1\); then
\(r_k\ge1\) would force \(\rho_k=1\), so eventually
\(r_k<1\) and \(\rho_k=r_k\).

Suppose a measurable estimator sequence \((\widehat V_k)\) has worst-case risks
\[
\mathcal W_k
=
\sup_{\mathbb P\in\mathcal D_{d_k,\epsilon_k}^{\mathrm{obs}}}
E_{\mathbb P}
[\{\widehat V_k-\Psi(\mathbb P)\}^2]
\longrightarrow0.
\]
The sample alphabet is finite. Thus each \(\widehat V_k\) is bounded on it, and these squared-error suprema are finite and nonnegative. The lower bound and the estimator infimum give, eventually,
\[
0\le\rho_k
\le\frac{\mathfrak R_{n_k,d_k,\epsilon_k}}{c_\star}
\le\frac{\mathcal W_k}{c_\star}.
\]
The squeeze theorem implies \(\rho_k\to0\), hence \(r_k\to0\).

Conversely, suppose \(r_k\to0\). Define a measurable estimator for every index by
\[
\widehat V_k=
\begin{cases}
\widehat V_{n_k,d_k,\epsilon_k}^{\mathrm{AF}},&n_k\ge1,\\
0,&n_k=0\text{ and }d_k<D_0,\\
1/2,&n_k=0\text{ and }d_k\ge D_0.
\end{cases}
\]
The zero-sample branches define constants on the singleton empty-sample space. For the eventual positive-sample indices, the established upper bound gives
\[
0\le
\sup_{\mathbb P\in\mathcal D_{d_k,\epsilon_k}^{\mathrm{obs}}}
E_{\mathbb P}
[\{\widehat V_k-\Psi(\mathbb P)\}^2]
\le C_\star\rho_k\longrightarrow0.
\]
A final application of the squeeze theorem proves uniform consistency and completes the equivalence.
% lean: CausalSmith.Stat.OptvalueVanishingoverlapRate.uniformConsistency_iff_of_frontier
\end{enumerate}
\end{proof}
